\documentclass[11pt]{article}

\usepackage[margin=2.5cm]{geometry}
\usepackage[utf8]{inputenc}
\usepackage[T1]{fontenc}
\usepackage{enumitem}
\usepackage{hyperref}
\usepackage{xcolor}
\usepackage{parskip}
\usepackage{titlesec}
\usepackage{fancyhdr}

\hypersetup{
  colorlinks=true,
  linkcolor=blue!60!black,
  urlcolor=blue!50!black,
  citecolor=blue!60!black
}

\titleformat{\section}{\Large\bfseries}{}{0em}{}
\titleformat{\subsection}{\large\bfseries}{}{0em}{}

\pagestyle{fancy}
\fancyhf{}
\rhead{\small EF PhD Fellowship 2026 --- Vincenzo Imperati}
\rfoot{\small\thepage}
\renewcommand{\headrulewidth}{0.4pt}

\title{\vspace{-1.5cm}\textbf{Verifiable Private Governance on Ethereum}\\[6pt]
\large Closing the Acceptance Gap for Private Collective Decisions\\[10pt]
\normalsize Ethereum Foundation PhD Fellowship 2026 --- Proposal\\
\normalsize RFP: Programmable Institutional Design and Verifiable Governance}

\author{Vincenzo Imperati\\
\small PhD Student in Computer Science, Sapienza University of Rome\\
\small \href{mailto:imperati@di.uniroma1.it}{imperati@di.uniroma1.it}
\; $\cdot$ \;
\href{https://vincenzo.imperati.dev}{vincenzo.imperati.dev}}
\date{}

\begin{document}
\maketitle
\thispagestyle{fancy}

%% ===================================================================
\section{Project Abstract}
%% ===================================================================

Ethereum can encode governance rules precisely, but private collective
decisions still often rely on trusted coordinators, tally operators, or
workflow controls that are hard to audit. This project studies a reusable
mechanism for closing the acceptance gap between an encrypted decision and the
proof that authorizes it, even when the encryption and proof systems are
cryptographically incompatible. The flagship case is private voting and
committee decision-making using threshold-Paillier for efficient encrypted
aggregation and Groth16 for proof verification that Ethereum can perform
cheaply on-chain. The fellowship would fund a stronger aggregate-integrity
mechanism built on top of that acceptance-layer base, a governance-level model
of what is immutable, adjustable, and discretionary, and a public artifact plus
open-access paper for Ethereum governance builders.


%% ===================================================================
\section{Objectives}
%% ===================================================================

\begin{enumerate}[leftmargin=*]
\item \textbf{Formalize a reusable mechanism for private collective
decisions.} Define clearly which institutional constraints become feasible
through programmability and verifiability in this setting, and map the design
into immutable, adjustable, and discretionary components, including who has
authority to set or change each part of the workflow.

\item \textbf{Strengthen tally-side integrity in a realistic way.} Starting from
the current transcript-auditable construction, design, implement, and formally
analyze an aggregate-integrity mechanism that reduces reliance on privileged
tally actors and closes the most important gap between accepted encrypted
ballots and the published aggregate.

\item \textbf{Produce reusable public infrastructure.} Release a public artifact
for Ethereum builders, including the core circuit and contract logic,
reproducible benchmarks, documentation, and a small reusable developer surface.

\item \textbf{Publish an open-access academic output.} Prepare and submit the
research to a named security venue target, with a public preprint path,
archived artifact, and permissive open-source licensing.
\end{enumerate}

\textbf{Success metrics.}
\begin{itemize}[leftmargin=*,nosep]
\item One paper submitted during the fellowship period, with a documented
      submission record and archived public artifact
\item One maintained public artifact release with documentation and reproducible results
\item One implemented and evaluated aggregate-integrity mechanism stronger than
      the current transcript-auditable baseline
\item One external round of feedback from Ethereum governance practitioners or
      adjacent projects
\end{itemize}


%% ===================================================================
\section{Outcomes}
%% ===================================================================

\subsection{From trusted procedure to enforceable rule}

The practical contribution is not ``blockchain voting'' in the abstract. It is
a reusable mechanism for higher-assurance private collective decisions on
Ethereum: committee votes, council approvals, board or membership decisions,
sealed governance proposals, and similar workflows where participants need
confidentiality without giving up public verifiability. The project targets
institutional and organizational settings, not national political elections
\cite{park2021}.

In a traditional confidential decision process, several roles are usually
separated: an administrator manages eligibility, trusted staff or scrutineers
handle ballot custody, and tally trustees or operators ensure that the object
counted is the same one originally cast. In many current Ethereum privacy
architectures, the same acceptance-side guarantee still falls to a coordinator
or application-specific off-chain procedure. The main institutional outcome of
this project is to move one specific rule from privileged trust into verifiable
code: once an encrypted decision is accepted, the object recorded on-chain must
be the same object authorized by the proof.

\subsection{A clearer allocation of authority}

This project speaks directly to the RFP's emphasis on predictable, enforceable,
and selectively controllable governance structures. In the governance design
studied here:

\begin{itemize}[leftmargin=*]
\item \textbf{Immutable:} the acceptance relation between the encrypted object
      and the proof that authorizes it; no coordinator, administrator, or tally
      actor should be able to override that relation after acceptance
\item \textbf{Adjustable:} governance parameters such as eligibility roots,
      thresholds, election timing, tally rules, and pause or restart policy,
      which remain under institutional control before or between governance
      rounds
\item \textbf{Discretionary:} optional modules above the cryptographic core,
      such as ballot-validity policies, decryption ceremonies, operational
      review, or institution-specific exception handling
\end{itemize}

This is intentionally a hybrid design rather than a fully automated one.
Programmability is used where discretionary override is most dangerous, while
institutions retain human control where policy, process, or legal context may
legitimately require it. The result is not only privacy, but a more legible
allocation of control: auditors and governance designers can see which rules are
hard-coded, which can be updated by authorized governance actors, and which are
left to explicit operational discretion.

\subsection{Real governance and compliance problems addressed}

The immediate problems addressed are not theoretical only. Many organizations
need procedures in which individual decisions remain confidential until a
process closes, while outsiders can still verify that only eligible inputs were
accepted and that the published aggregate corresponds to those accepted inputs.
This arises in grant committees, professional associations, protocol councils,
university bodies, and other collectives that need confidentiality without
opaque administration. The project therefore treats auditable governance and
verifiable compliance as core features rather than side benefits.

\subsection{Benefit to the wider Ethereum ecosystem}

Ethereum has strong governance tooling, but privacy-preserving governance still
often relies on specialized workflow actors, off-chain coordination, or design
tradeoffs that are difficult to reuse. This project benefits the ecosystem in
three ways:

\begin{itemize}[leftmargin=*]
\item it turns a currently under-specified mixed-domain composition problem into
      a documented and reusable design pattern;
\item it gives Ethereum builders a public artifact and benchmark baseline for a
      governance problem that sits directly at the intersection of encryption,
      succinct proofs, and EVM constraints shaped by Ethereum's BN254 precompiles
      \cite{eip196,eip197};
\item it translates a narrow cryptographic contribution into institutional
      design language that can be understood by non-cryptographers working on
      governance products and frameworks.
\end{itemize}


%% ===================================================================
\section{Grant Scope}
%% ===================================================================

\subsection{Current starting point}

This proposal does not start from a blank slate. I already have a public
research artifact centered on the acceptance side of the problem, including the
acceptance gap between encrypted decisions and their on-chain authorization
proof: a Circom circuit, a Solidity contract using a real Groth16 verifier on
EVM, a reproduced attack on the naive composition, passing tests, and measured
benchmarks at realistic parameters. I also have a paper draft and a detailed
map of claims, evidence, and open gaps.

This existing base reduces execution risk, but the work is not already
complete. The current artifact is strongest on acceptance-side integrity and on
public auditability of the tally transcript. The fellowship would fund the next
substantive step: turning that base into a clearer governance contribution with
stronger tally-side guarantees, public packaging, and a publication-ready
research output.

\subsection{Research work funded by the fellowship}

\textbf{Work Package 1: Formal model and governance framing.}
I will refine the current technical model into a governance-oriented formal
description of the mechanism, including:

\begin{itemize}[leftmargin=*]
\item the mixed-domain threat model;
\item the mapping between cryptographic guarantees and institutional
      constraints;
\item a clear immutable / adjustable / discretionary decomposition with a
      specified authority model; and
\item a comparison with both traditional trustee- or coordinator-run workflows
      and nearby Ethereum governance architectures.
\end{itemize}

\textbf{Work Package 2: Aggregate-integrity mechanism.}
I will design, implement, and evaluate a mechanism that strengthens the current
tally path beyond mismatch detectability. The primary path is an incremental
aggregate-integrity construction in which the tally-side aggregate becomes a
deterministic function of accepted ciphertexts by construction, reducing room
for operator discretion after acceptance. If that primary path proves too costly
or too weak for a precise claim, the fallback is a narrower tally-consistency
upgrade that still materially improves verifiability over the present baseline.

Any retained construction must satisfy four conditions:

\begin{itemize}[leftmargin=*]
\item it yields a precise formal claim;
\item it has clear implementation semantics;
\item it can be tested and benchmarked convincingly;
\item it does not distort the project away from its governance-mechanism core.
\end{itemize}

\textbf{Work Package 3: Public artifact and ecosystem packaging.}
I will mature the existing public repository into a reusable artifact release
whose packaging is centered on the artifact itself rather than on the broader
research history behind it. This will include:

\begin{itemize}[leftmargin=*]
\item the core circuit and contract logic;
\item benchmark data and regeneration instructions;
\item threat-model and limitation notes;
\item developer-facing documentation; and
\item a small reusable package or integration surface for Ethereum builders.
\end{itemize}

\textbf{Work Package 4: Academic output and dissemination.}
I will write and submit the paper, target IEEE S\&P 2027 Cycle 2 as the
primary venue and ACM CCS 2027 as fallback, archive the artifact, and produce a
concise non-technical explanation for Ethereum governance practitioners.

\subsection{Expected outputs}

\begin{itemize}[leftmargin=*,nosep]
\item One open-access research paper submitted to IEEE S\&P 2027 Cycle 2 or ACM CCS 2027
\item One public preprint or equivalent open-access manuscript path
\item One documented public open-source artifact release under a permissive license
\item One archived artifact snapshot, for example Zenodo
\item One concise practitioner-facing technical explainer for the Ethereum
      governance ecosystem
\end{itemize}


%% ===================================================================
\section{Related Work}
%% ===================================================================

This project sits between institutional governance practice and cryptographic
systems work that rarely meet cleanly in Ethereum deployments.

On the governance side, confidential decision workflows traditionally rely on a
separation of duties among administrators, scrutineers, or tally trustees.
Ethereum already has mature systems for collective decision-making, ranging from
off-chain signaling to on-chain execution frameworks, but private governance
architectures encode or retain those roles in different ways. MACI shows that
private on-chain governance is valuable and feasible, but its audited
architecture still organizes tallying and proof generation around a coordinator
role. Open Vote Network shows a different design point through self-tallying
voting, avoiding the threshold-encrypted aggregation path altogether. Helios and
the follow-up critique literature remain important reminders that
composition-level gaps, not only broken primitives, can undermine governance
systems \cite{maci2024audit,mccorry2017,adida2008,cortier2013,bernhard2015}.

On the cryptographic side, same-domain composition is much better understood.
LegoSNARK provides a powerful modular composition framework inside a
proof-friendly domain. SAVER and newer encryption-friendly proof systems such as
Tangram illustrate how much easier composition becomes when encrypted objects
and proof statements remain within compatible algebraic structures. Recent work
on Paillier proofs, including efficient zero-knowledge arguments and direct
range proofs, shows that Paillier-side reasoning is possible, but also that it
is not free in practice. Modern threshold Paillier systems such as Tiresias
show that Paillier remains a serious primitive rather than only a historical
one \cite{legosnark2019,saver2019,tangram2025,gong2024,xie2024,tiresias2025}.

This broader relevance is not unique to voting. In federated learning, secure
aggregation, verifiable aggregation, and input verification are active problems,
which reinforces the point that encrypted aggregation plus public or
third-party-checkable correctness is an important systems pattern in its own
right \cite{veriflfix2022,lzksa2025}. I am not proposing a full federated
learning deliverable in this fellowship, but the literature helps show that the
composition pattern isolated here is not an eccentric one-off artifact of
voting.

The specific gap addressed here is narrower and more practical: Ethereum still
lacks a clean, reusable mechanism for the case where a raw encrypted governance
object must be authorized by a proof that Ethereum can verify cheaply on-chain,
despite the encryption and proof systems living in incompatible cryptographic
domains. That is the acceptance gap I aim to formalize and close. Existing systems either remain in same-domain
cryptographic settings, avoid threshold-encrypted aggregation altogether, or
accept different workflow assumptions. The project therefore does not claim that
encryption, succinct proofs, or private governance are individually missing; it
claims that their mixed-domain interface on Ethereum remains under-specified and
under-packaged.


%% ===================================================================
\section{About the Applicant}
%% ===================================================================

I am a PhD student in Computer Science at Sapienza University of Rome, where I
work under the supervision of Alessandro Mei. My broader PhD work is in
security, large-scale systems, and open protocols, and I am currently a
Visiting Researcher at the University of Zurich. I have published at
\textit{USENIX Security} and \textit{IEEE Access}, and I have spent the last
years building public infrastructure as well as doing academic research. This
proposal sits inside my PhD research environment rather than outside it: the
fellowship would fund a line of work already grounded in university research,
existing supervision, and a working public artifact.

I am well suited to this fellowship because this project requires exactly the
combination I already have: security-research discipline, real artifact
engineering, and practical Ethereum experience. I founded and maintain
\textit{BigBrotr}, an OpenSats-funded observability platform for the Nostr
protocol, which demonstrates my ability to deliver funded public research
infrastructure. I have also built multiple EVM projects, including MeltyFi,
Nostreum, and the earlier on-chain election-integrity prototype from which this
research direction emerged. This means I am not approaching Ethereum governance
as an external observer; I am already an active builder in the ecosystem.

\textbf{Relevant links}
\begin{itemize}[leftmargin=*,nosep]
\item Personal website: \href{https://vincenzo.imperati.dev}{vincenzo.imperati.dev}
\item Sapienza researcher profile: \href{https://research.uniroma1.it/researcher/0f9c3c1227f530e4af686c68b186e30bf4f6679206966e380af17780}{research.uniroma1.it}
\item GitHub: \href{https://github.com/VincenzoImp}{github.com/VincenzoImp}
\item USENIX Security paper: \href{https://www.usenix.org/conference/usenixsecurity25/presentation/imperati}{The Conspiracy Money Machine}
\item IEEE Access paper: \href{https://ieeexplore.ieee.org/document/11224376}{Hard to See, Harder to Block}
\item Project repository link: \href{https://github.com/VincenzoImp/paillier-groth16-binding}{github.com/VincenzoImp/paillier-groth16-binding}
\item ORCID: \href{https://orcid.org/0009-0001-9437-1384}{0009-0001-9437-1384}
\item Google Scholar: \href{https://scholar.google.com/citations?user=V4XgGVcAAAAJ}{scholar.google.com/citations?user=V4XgGVcAAAAJ}
\item Supervisor academic profile: \href{http://wwwusers.di.uniroma1.it/~mei/}{Alessandro Mei}
\end{itemize}


%% ===================================================================
\section{Methodology}
%% ===================================================================

This is applied cryptography research with a concrete institutional example.
The methodological goal is not to simulate a full voting platform; it is to
study, strengthen, and package one reusable governance mechanism that becomes
useful once private decisions and succinct authorization proofs must coexist on
Ethereum.

\subsection{Phase 1: Formalization and architecture comparison}

I will start by refining the mixed-domain model from the current research base
into a cleaner governance-level statement. This includes:

\begin{itemize}[leftmargin=*]
\item specifying the threat model induced by a public bulletin board and cheap
      on-chain proof verification;
\item mapping the mechanism into immutable, adjustable, and discretionary
      governance components together with the authority structure around each;
\item comparing it against adjacent design points, especially
      coordinator-centric, self-tallying, and same-domain approaches; and
\item making explicit the comparison with traditional trustee- or
      scrutineer-based governance workflows.
\end{itemize}

\subsection{Phase 2: Design and evaluation of the aggregate-integrity upgrade}

I will then build the main research extension of the fellowship: a stronger
aggregate-integrity mechanism. I already know what the present baseline is:

\begin{itemize}[leftmargin=*]
\item accepted ballots are bound correctly at the board layer;
\item transcript-derived aggregate mismatch is publicly detectable.
\end{itemize}

The fellowship will fund the step beyond that baseline. I will first implement
the primary path described in the Grant Scope: an incremental
aggregate-integrity mechanism that makes the published aggregate follow from the
accepted ciphertext set more directly. If that path does not support a clear and
credible claim after implementation and measurement, I will fall back to a
narrower tally-consistency upgrade rather than forcing a weak result into the
paper.

\subsection{Phase 3: Public artifact, documentation, and paper}

Once the research line is stable, I will:

\begin{itemize}[leftmargin=*]
\item stabilize the existing public repository;
\item write a clear threat-model and limitations document;
\item archive benchmark outputs and regeneration paths;
\item prepare the artifact for archival release; and
\item produce a short explainer for practitioners who care about governance
      design but are not specialists in cryptography.
\end{itemize}

\subsection{Phase 4: Submission and dissemination}

In the final phase, I will turn the research into a formal academic output and a
usable ecosystem deliverable. This includes rewriting the paper into
venue-ready form, submitting it to the target venue, preserving a public
manuscript path consistent with venue policy, and collecting external feedback
from adjacent Ethereum governance builders.

\subsection{Hybrid design and scope discipline}

The project is intentionally hybrid rather than fully automated. Acceptance-side
binding and the corresponding aggregate-integrity mechanism are the parts that
should become maximally rigid and auditable. Decryption ceremonies,
institution-specific ballot policies, and exceptional operational decisions may
remain partially off-chain or discretionary, because the purpose of the project
is to clarify the boundary between programmable enforcement and retained human
control, not to eliminate institutions from the loop.

The project will remain intentionally scoped. During the fellowship I will not
reframe it as:

\begin{itemize}[leftmargin=*]
\item a complete national-election system;
\item a complete solution to all tally and decryption correctness questions;
\item or a three-application general framework spanning voting, federated
      learning, and MPC simultaneously.
\end{itemize}

This scope discipline is part of the execution strategy: the goal is to deliver
one strong, useful result rather than many half-finished ones.


%% ===================================================================
\section{Timeline}
%% ===================================================================

Planned fellowship period:

\begin{itemize}[leftmargin=*,nosep]
\item June 2026 to May 2027
\item total requested fellowship support: \$24{,}000
\item approximate total effort used for planning: 720 hours
\item indicative planning rate used for milestone allocation: approximately \$33.33/hour
\item milestone budgets below distribute the fixed fellowship amount across the planned work packages
\end{itemize}

\subsection*{Milestone 1: Institutional model and architecture comparison}

\textbf{Timing:} Months 1--3\\
\textbf{Budget:} \$6{,}000\\
\textbf{Number of hours (roughly):} 180

\textbf{Work}
\begin{itemize}[leftmargin=*,nosep]
\item consolidate the formal and governance model
\item finalize the immutable / adjustable / discretionary decomposition
\item compare the mechanism against trustee-run, coordinator-run, and adjacent
      Ethereum governance designs
\end{itemize}

\textbf{Deliverables}
\begin{itemize}[leftmargin=*,nosep]
\item stable problem statement
\item governance-comparison section including traditional workflow counterpart
\item updated research plan for the aggregate-integrity extension
\end{itemize}

\subsection*{Milestone 2: Aggregate-integrity mechanism}

\textbf{Timing:} Months 3--6\\
\textbf{Budget:} \$8{,}000\\
\textbf{Number of hours (roughly):} 240

\textbf{Work}
\begin{itemize}[leftmargin=*,nosep]
\item implement the primary aggregate-integrity path
\item implement the mechanism in the existing artifact stack
\item add tests and measure cost
\item formalize the corresponding claim, or narrow the result to the strongest
      defensible fallback
\end{itemize}

\textbf{Deliverables}
\begin{itemize}[leftmargin=*,nosep]
\item implemented and benchmarked aggregate-integrity upgrade
\item updated threat model and limitation notes
\item locked final paper scope
\end{itemize}

\subsection*{Milestone 3: Public artifact stabilization and archival release}

\textbf{Timing:} Months 6--8\\
\textbf{Budget:} \$5{,}000\\
\textbf{Number of hours (roughly):} 150

\textbf{Work}
\begin{itemize}[leftmargin=*,nosep]
\item stabilize the existing public repository
\item write documentation and reproduction instructions
\item package benchmark data and artifact metadata
\item archive a stable tagged snapshot
\end{itemize}

\textbf{Deliverables}
\begin{itemize}[leftmargin=*,nosep]
\item documented public artifact release
\item archived release snapshot
\item practitioner-facing explainer draft
\end{itemize}

\subsection*{Milestone 4: Paper submission and dissemination}

\textbf{Timing:} Months 8--12\\
\textbf{Budget:} \$5{,}000\\
\textbf{Number of hours (roughly):} 150

\textbf{Work}
\begin{itemize}[leftmargin=*,nosep]
\item complete the paper
\item submit to IEEE S\&P 2027 Cycle 2 or ACM CCS 2027 fallback
\item collect ecosystem feedback
\item publish the public explainer and integration notes
\end{itemize}

\textbf{Deliverables}
\begin{itemize}[leftmargin=*,nosep]
\item submitted paper and submission record
\item public preprint path or submission-ready manuscript
\item final public explainer
\item short feedback memo from ecosystem conversations
\end{itemize}


\begin{thebibliography}{99}\small

\bibitem{eip196}
Christian Reitwiessner.
\newblock {EIP-196}: Precompiled contracts for addition and scalar multiplication on the elliptic curve alt\_bn128.
\newblock Ethereum Improvement Proposals, no.~196, 2017.
\newblock \url{https://eips.ethereum.org/EIPS/eip-196}

\bibitem{eip197}
Vitalik Buterin and Christian Reitwiessner.
\newblock {EIP-197}: Precompiled contracts for optimal ate pairing check on the elliptic curve alt\_bn128.
\newblock Ethereum Improvement Proposals, no.~197, 2017.
\newblock \url{https://eips.ethereum.org/EIPS/eip-197}

\bibitem{maci2024audit}
Kyle Charbonnet and Yufei Li.
\newblock {MACI} Security Audit.
\newblock PSE Security, July 2024.
\newblock \url{https://maci.pse.dev/assets/files/20240731_PSE_Audit_audit_report-a0a0f08a5c621ccd0389d5e345c119be.pdf}

\bibitem{mccorry2017}
Patrick McCorry, Siamak F. Shahandashti, and Feng Hao.
\newblock A smart contract for boardroom voting with maximum voter privacy.
\newblock In {\em Financial Cryptography and Data Security}, pages 357--375, 2017.
\newblock doi: \url{https://doi.org/10.1007/978-3-319-70972-7_20}

\bibitem{adida2008}
Ben Adida.
\newblock Helios: Web-based Open-Audit Voting.
\newblock In {\em 17th USENIX Security Symposium}, pages 335--348, 2008.
\newblock \url{https://www.usenix.org/conference/17th-usenix-security-symposium/helios-web-based-open-audit-voting}

\bibitem{cortier2013}
V{\'e}ronique Cortier and Ben Smyth.
\newblock Attacking and fixing Helios: An analysis of ballot secrecy.
\newblock {\em Journal of Computer Security}, 21(1):89--148, 2013.

\bibitem{bernhard2015}
David Bernhard, V{\'e}ronique Cortier, David Galindo, Olivier Pereira, and Bogdan Warinschi.
\newblock {SoK}: A comprehensive analysis of game-based ballot privacy definitions.
\newblock In {\em IEEE Symposium on Security and Privacy}, pages 499--516, 2015.
\newblock doi: \url{https://doi.org/10.1109/SP.2015.37}

\bibitem{park2021}
Sunoo Park, Michael Specter, Neha Narula, and Ronald L. Rivest.
\newblock Going from bad to worse: from Internet voting to blockchain voting.
\newblock {\em Journal of Cybersecurity}, 7(1):tyaa025, 2021.
\newblock doi: \url{https://doi.org/10.1093/cybsec/tyaa025}

\bibitem{legosnark2019}
Matteo Campanelli, Dario Fiore, and Ana{\"i}s Querol.
\newblock LegoSNARK: Modular design and composition of succinct zero-knowledge proofs.
\newblock In {\em ACM CCS}, pages 2075--2092, 2019.
\newblock doi: \url{https://doi.org/10.1145/3319535.3339820}

\bibitem{saver2019}
Jiwon Lee, Jaekyoung Choi, Jihye Kim, and Hyunok Oh.
\newblock SAVER: SNARK-friendly, additively homomorphic, and verifiable encryption and decryption with rerandomization.
\newblock Cryptology ePrint Archive, Paper 2019/1270.

\bibitem{tangram2025}
Gweonho Jeong, Myeongkyun Moon, Geonho Yoon, Hyunok Oh, and Jihye Kim.
\newblock Tangram: Encryption-friendly SNARK framework under Pedersen committed engines.
\newblock Cryptology ePrint Archive, Paper 2025/540.

\bibitem{gong2024}
Borui Gong, Wang Fat Lau, Man Ho Au, Rupeng Yang, Haiyang Xue, and Lichun Li.
\newblock Efficient Zero-Knowledge Arguments for Paillier Cryptosystem.
\newblock Cryptology ePrint Archive, Paper 2024/1303.

\bibitem{xie2024}
Zhikang Xie, Mengling Liu, Haiyang Xue, Man Ho Au, Robert H. Deng, and Siu-Ming Yiu.
\newblock Direct Range Proofs for Paillier Cryptosystem and Their Applications.
\newblock Cryptology ePrint Archive, Paper 2024/1355.

\bibitem{tiresias2025}
Offir Friedman, Avichai Marmor, Dolev Mutzari, Yehonatan C. Scaly, Yuval Spiizer, and Avishay Yanai.
\newblock Tiresias: Large Scale, Maliciously Secure Threshold Paillier.
\newblock Cryptology ePrint Archive, Paper 2023/998.

\bibitem{veriflfix2022}
Xiaojie Guo.
\newblock Fixing Issues and Achieving Maliciously Secure Verifiable Aggregation in ``VeriFL: Communication-Efficient and Fast Verifiable Aggregation for Federated Learning''.
\newblock Cryptology ePrint Archive, Paper 2022/1073.

\bibitem{lzksa2025}
Zhi Lu and Songfeng Lu.
\newblock LZKSA: Lattice-based special zero-knowledge proofs for secure aggregation's input verification.
\newblock Cryptology ePrint Archive, Paper 2025/1141.

\end{thebibliography}

\end{document}
